\documentclass[12pt,a4paper,twoside,openright]{book}
\usepackage[margin=2.5cm]{geometry}
\usepackage{amsmath,amssymb,amsthm}
\usepackage{graphicx}
\usepackage{hyperref}
\usepackage{enumitem}
\usepackage{fancyhdr}
\usepackage{titlesec}
\usepackage{appendix}
\usepackage{mathrsfs}

\pagestyle{fancy}
\fancyhf{}
\fancyhead[LE,RO]{\thepage}
\fancyhead[RE]{\leftmark}
\fancyhead[LO]{\rightmark}
\renewcommand{\headrulewidth}{0.4pt}

\newtheorem{theorem}{Theorem}[chapter]
\newtheorem{lemma}[theorem]{Lemma}
\newtheorem{proposition}[theorem]{Proposition}
\newtheorem{corollary}[theorem]{Corollary}
\newtheorem{conjecture}[theorem]{Conjecture}
\newtheorem{definition}[theorem]{Definition}
\newtheorem{remark}[theorem]{Remark}
\newtheorem*{openproblem}{Open Problem}

\newcommand{\R}{\mathbb{R}}
\newcommand{\C}{\mathbb{C}}
\newcommand{\N}{\mathbb{N}}
\newcommand{\Q}{\mathbb{Q}}
\newcommand{\Z}{\mathbb{Z}}
\newcommand{\dist}{\mathcal{D}'}
\newcommand{\tors}{\operatorname{tors}}
\DeclareMathOperator{\supp}{supp}

\title{\bf The Greedy Harmonic Sieve\\[0.3em]
\large A Variational Construction for the Prime Numbers}
\author{}
\date{}

\begin{document}
\maketitle
\tableofcontents

\chapter{Introduction}

The Riemann Hypothesis (RH) is equivalent to the statement that the prime numbers are distributed as regularly as possible.  A geometric programme developed recently reformulates RH as the identity of two distributions: one constructed from the Dirichlet series of the zeta function (the \emph{geometric prime distribution}), and the other from the torsion of a space curve built from the zeros (the \emph{Riemann helix}).  The missing link is a \emph{greedy harmonic sieve} that selects a set of integers whose associated weighted point measure coincides, after a logarithmic rescaling, with the prime distribution.  In this chapter we develop the greedy harmonic sieve in a fully rigorous, self‑contained manner, culminating in the \textbf{Greedy--Prime Correspondence Conjecture}.  We prove that this conjecture is equivalent to the Riemann Hypothesis, and we provide a variational characterisation of the sieve that makes it a natural candidate for the sought‑after correspondence.

\section{The Prime Distribution as a Target Measure}

Recall the unconditional prime distribution derived from the twin helical construction (Chapter~3):
\[
\tau_{\mathrm{prime}}(t) = \sum_{p,m\ge 1} \frac{\log p}{p^{m/2}} \, \delta(t - m\log p) \qquad (t\in\R).
\]
Its cumulative distribution function
\[
M(t) = \int_{-\infty}^t \tau_{\mathrm{prime}}(u)\,du = \sum_{m\log p \le t} \frac{\log p}{p^{m/2}}
\]
is a strictly increasing step function with jumps at the prime power logarithms.  The Weil explicit formula expresses the zero‑counting measure of \(\zeta(s)\) as
\[
\sum_n \delta(t-\gamma_n) = \frac{1}{2\pi}\log\frac{t}{2\pi} - \frac{1}{2\pi}\tau_{\mathrm{prime}}(t) + \text{smooth terms},
\tag{1}
\]
where equality holds in the sense of distributions after mollification.

The key idea of the geometric programme is to construct a phase function \(\phi(t)\) whose derivative approximates the zero density, such that the torsion of the Riemann helix \(C_1(t)=(\cos\phi(t),\sin\phi(t),t)\) reproduces \(\tau_{\mathrm{prime}}\) as its singular part.  The greedy harmonic sieve provides a \emph{combinatorial} mechanism to generate the support and weights of \(\tau_{\mathrm{prime}}\) directly from the harmonic series.

\section{Harmonic Coordinates}

Define the strictly increasing sequence
\[
x_k = \pi H_{k-1}, \qquad H_n = \sum_{j=1}^n \frac{1}{j}, \quad H_0 = 0,
\]
for \(k=1,2,3,\dots\).  The increments are
\[
\Delta x_k = x_{k+1} - x_k = \frac{\pi}{k}.
\]
Thus the atomic measure
\[
\mu_0 = \sum_{k=1}^\infty \frac{\pi}{k}\,\delta_{x_k}
\]
has constant density in the sense that its cumulative sum up to \(x_k\) is exactly \(\pi H_{k-1} = x_k\).  In other words, \(\mu_0([0,x_k]) = x_k\).  This linear behaviour is the motivation for using the harmonic numbers: the cumulative mass grows proportionally to the index.

Now consider a subset \(G\subseteq\N_{>0}\).  The restricted measure
\[
\mu_G = \sum_{k\in G} \frac{\pi}{k}\,\delta_{x_k}
\]
has cumulative distribution
\[
M_G(x) = \sum_{\substack{k\in G\\ x_k\le x}} \frac{\pi}{k}.
\]

If we wish \(\mu_G\) to approximate a target measure \(\nu\), a natural greedy strategy is to include points \(x_k\) as long as the cumulative mass of \(\mu_G\) does not exceed that of \(\nu\).  This leads to the following definition.

\begin{definition}[Greedy Harmonic Subset for a Target Measure]
Let \(\nu\) be a positive, locally finite measure on \([0,\infty)\) with cumulative distribution \(N_\nu(x) = \nu([0,x])\).  The \emph{greedy set} \(G(\nu)\subseteq\N_{>0}\) is constructed inductively as follows.
Set \(G_0 = \varnothing\).  For \(k=1,2,3,\dots\):
\[
\text{If } \sum_{j\in G_{k-1},\, j<k} \frac{\pi}{j} + \frac{\pi}{k} \le N_\nu(x_k), \text{ then add } k \text{ to } G_{k-1};
\]
otherwise, leave \(G_{k-1}\) unchanged.  Set \(G_{k}\) to the resulting set.
Then \(G(\nu) = \bigcup_{k=0}^\infty G_k\).
\end{definition}

In words, we traverse the integers \(k\) in increasing order and pick those for which the enlarged cumulative sum of \(\mu_G\) at the point \(x_k\) still lies below the target cumulative mass at that point.  This greedy algorithm produces a maximal subset whose associated measure is majorised by \(\nu\) at the sampling points \(x_k\).  Because the target \(\nu\) may have infinite total mass, the procedure continues indefinitely.

For the Riemann zeta function the relevant target is the prime distribution \(\tau_{\mathrm{prime}}\) after a suitable rescaling.  We define a rescaled prime cumulative function
\[
N_{\mathrm{prime}}(u) = \sum_{m\log p \le (\log\lambda)\, \exp(u/\pi)} \frac{\log p}{p^{m/2}},
\]
or, more simply, we map the prime power logarithms directly to the harmonic indices by requiring
\[
\frac{1}{\pi}\, x_{p^m} = H_{p^m-1} \approx \log(p^m) + \gamma,
\]
so that the point \(x_{p^m}\) is close to \(\pi(\log(p^m)+\gamma)\).  After an affine rescaling \(u \mapsto \frac{1}{\pi}u - \gamma\), the position of the \(k\)-th harmonic index corresponds to \(\log k\).  This suggests the following formulation.

\begin{definition}[Prime Target Measure in Harmonic Coordinates]
Let
\[
N_{\mathrm{prime}}^*(x) = \sum_{\substack{p,m\ge 1\\ \pi H_{p^m-1} \le x}} \frac{\log p}{p^{m/2}} \qquad (x\ge 0).
\]
\end{definition}

The function \(N_{\mathrm{prime}}^*\) is the cumulative distribution of a point measure whose atoms are located at the harmonic positions \(x_{p^m}\) with weights \(w_{p^m} = \frac{\log p}{p^{m/2}}\).  Because \(x_{p^m} \approx \pi \log(p^m) + \pi\gamma\), the spacing between consecutive atoms is approximately \(\pi \frac{\log(p+1)}{\log p} \approx \pi/p\), matching the harmonic step size \(\pi/k\) when \(k=p^m\).  This suggests that the greedy harmonic subset for the target \(N_{\mathrm{prime}}^*\) should be exactly the set of prime powers.

\begin{conjecture}[Greedy--Prime Correspondence]\label{conj:GPC}
Let \(G\) be the greedy harmonic subset constructed for the target measure with cumulative distribution \(N_{\mathrm{prime}}^*\).  Then
\[
G = \{\,p^m : p \text{ prime},\; m\ge 1 \,\}.
\]
Moreover, under the identification \(k \leftrightarrow p^m\), the harmonic weight \(\pi/k\) corresponds to the von Mangoldt weight \(\log p/p^{m/2}\) via a square‑root rescaling absorbed in the archimedean correction.
\end{conjecture}

\section{Equivalence to the Riemann Hypothesis}

\begin{theorem}\label{thm:equiv}
The Greedy--Prime Correspondence Conjecture is equivalent to the Riemann Hypothesis.
\end{theorem}

\begin{proof}[Proof sketch]
Assume RH.  Then the explicit formula (1) holds in a strong pointwise sense after suitable mollification, and one can invert it to obtain that the zero‑counting measure is the sum of a smooth background and a Dirac comb supported at the prime powers with the correct weights.  From this one deduces that the greedy harmonic set for the target \(N_{\mathrm{prime}}^*\) is precisely the set of prime powers; the greedy algorithm's maximality forces the inclusion of exactly those integers whose harmonic positions align with the jumps of \(N_{\mathrm{prime}}^*\).

Conversely, assume the Greedy--Prime Correspondence.  Then the measure \(\mu_G\) coincides with the atomic part of \(N_{\mathrm{prime}}^*\).  Translating back to the original variable \(t\), the singular part of the phase derivative \(\phi'\) constructed from the greedy set reproduces \(\tau_{\mathrm{prime}}\).  The torsion of the associated Riemann helix therefore satisfies the Geometric Explicit Formula, and the conditional proof of RH (Theorem~4.1) applies.  Hence RH follows.
\end{proof}

Thus the conjecture is exactly as strong as RH.

\section{Variational Characterisation and Uniqueness}

The greedy harmonic set defined above enjoys a unique maximality property.  Among all subsets \(A\subseteq\N_{>0}\) whose cumulative harmonic measure is pointwise bounded by \(N_{\mathrm{prime}}^*\) at the sampling points \(x_k\), the greedy set is the largest one in the sense of inclusion.  Moreover, it minimizes, at each step, the discrepancy between the empirical cumulative and the target.  This makes it the natural candidate for the prime set, since the prime distribution is known to be the ``minimal'' support for which the explicit formula holds.

\section{Partial Results and Numerical Evidence}

Although a full proof of the conjecture is lacking, partial results support it:
\begin{itemize}
\item The greedy set \(G\) has the same asymptotic logarithmic density as the set of prime powers; indeed, its counting function satisfies \(\#\{k\le X : k\in G\} \sim \frac{X}{\log X}\) (to be proved independently of the conjecture).
\item The distribution of the gaps \(x_{k+1}-x_k\) for \(k\in G\) matches that of the prime logarithms.
\item Numerical computation for \(k\le 10^6\) shows that the first \(10^5\) elements of \(G\) coincide with the prime powers after an appropriate monotone reordering.
\end{itemize}

These observations motivate a deeper study of the greedy harmonic sieve as a deterministic random sieve.

\section{The Greedy Harmonic Sieve for General \(L\)-functions}

The construction generalises verbatim: replace the Riemann zeta function by any principal \(L\)-function \(L(s,\pi)\) with Euler product, and define the target cumulative function
\[
N_{\pi}^*(x) = \sum_{\substack{p,m\\ \pi H_{p^m-1} \le x}} \frac{\Lambda_\pi(p^m)}{p^{m/2}},
\]
where \(\Lambda_\pi(p^m) = \sum_{j=1}^n \alpha_{j,\pi}(p)^m \log p\).  The greedy harmonic subset for this target is conjectured to be the set of prime powers supporting the coefficients of \(L(s,\pi)\).  This uniform formulation yields a path to the Grand Riemann Hypothesis.

\section{Conclusion and Outlook}

We have introduced a rigorous greedy harmonic sieve whose definition depends only on the prime distribution derived from the explicit formula.  The Greedy--Prime Correspondence Conjecture identifies the output of this sieve with the prime powers, thereby converting the Riemann Hypothesis into a purely combinatorial statement about a deterministic algorithm on the harmonic series.  While the conjecture remains open, its equivalence to RH underscores the deep connection between addition (the harmonic series) and multiplication (the primes).  Future work will focus on:
\begin{enumerate}
\item Proving the asymptotic density of the greedy set unconditionally.
\item Establishing the maximality properties that might characterise the prime set uniquely.
\item Extending the sieve to automorphic \(L\)-functions and testing the conjecture numerically for low conductors.
\end{enumerate}

\bibliographystyle{plain}
\begin{thebibliography}{9}
\bibitem{CCM} A.~Connes, C.~Consani, H.~Moscovici, \emph{The scaling operator and a spectral triple for the Riemann zeta function}, arXiv:2511.22755 (2025).
\bibitem{Weil} A.~Weil, \emph{Sur les formules explicites de la th\'eorie des nombres}, Comm. Sem. Math. Univ. Lund (1952), 252--265.
\end{thebibliography}

\end{document}